\ch{Chapter 5: Conclusions and Future Research}\label{chap:Chapter_5}

This dissertation set out to learn whether Information Technology and Information Assurance professionals are aware of Digital Engineering capabilities and whether they perceive those capabilities as valuable for their work. Chapter~1 framed the problem and posed the three research questions. Chapter~2 exposed a near-complete absence of academic research carrying Digital Engineering into enterprise IT and Information Assurance, and it argued that adoption in these domains is governed by perception rather than by proven technical merit. Chapter~3 specified a quantitative, cross-sectional survey methodology, and Chapter~4 reported its results from a sample of 1{,}411 professionals. The present chapter concludes the study. It synthesizes those results, interprets them against the literature, states the conclusions warranted for each research question, revisits the work's significance along the dimensions established in Chapter~1, acknowledges the limitations that bound the conclusions, and charts future research.

\section{Summary of the Study}

Evidence came from a quantitative survey of IT and Information Assurance professionals, structured around six thematic sections that measured awareness, recognition of specific capabilities, perceived applicability, and value assessments for both domains. Within that instrument, explicit traceability connected survey questions to the three research questions. Analytical plans specified descriptive characterization, composite construction with reliability assessment, comparative testing across professional field and experience level, association analysis, and two planned bias diagnostics. A realized sample of 1{,}411 completed responses exceeded the planned target of 385 to 450 by roughly threefold. This provided high precision while shifting emphasis away from statistical significance, which is near-universal at this sample size, and toward effect sizes and confidence intervals.

\section{Conclusions by Research Question}

\subsection{Research Question 1: Awareness}

To what extent are IT and Information Assurance professionals aware of Digital Engineering capabilities? The evidence yields a two-part answer. Awareness of Digital Engineering as a named discipline is low: 78.2\% of respondents placed themselves in the two lowest familiarity categories, the familiarity item mean was 1.94, and only 22.2\% reported any professional exposure within the prior two years. Endorsement of the specific capabilities that constitute Digital Engineering was nonetheless near-universal once the instrument described those capabilities in functional language, with 93.9\% agreement across the Section~2 capability-recognition items. This divergence reads most defensibly as low awareness of the umbrella term coexisting with high endorsement of the underlying capability descriptions.

Two bounds attach to that second half. Section~2 presented functional descriptions and asked for agreement, so it measures recognition of a capability once described and leaves prior knowledge of Model-Based Systems Engineering, digital threads, digital twins, and Product Lifecycle Management unmeasured. Acquiescence to affirmatively framed statements would produce the same near-absence of disagreement. Research Question~1 therefore warrants a narrow conclusion: awareness of Digital Engineering as a named discipline is low, and the described capabilities draw near-universal endorsement when put to respondents in plain functional language. The finding echoes the awareness deficit that \citeauthor{Henderson_McDermott_Salado_2024} documented even within the systems engineering discipline, where roughly 22\% of practitioners could not articulate a clear definition of Model-Based Systems Engineering~\parencite{Henderson_McDermott_Salado_2024}. If awareness is fragile within the originating discipline, the low familiarity observed here among professionals who have never systematically encountered these methods is unsurprising and substantively important.

\subsection{Research Question 2: Perceived Value}

Research Question~2 asks to what extent professionals perceive these capabilities as valuable. The evidence admits a clear answer. The strength of that answer lies in the precision of the estimates rather than in the difficulty of disconfirming the affirmatively framed items, a limit Chapter~4 flags. Perceived value was high and precisely estimated across the applicability (86.3\% agreement), IT-value (87.7\%), and security-value (82.2\%) blocks, with narrow Wilson confidence intervals reflecting the large sample. Interest and intent corroborated the value findings, as 89.0\% expressed interest in learning more and 95.9\% indicated they would recommend their organization explore adoption. The two investment-willingness items registered near-unanimous willingness conditional on demonstrated return on investment, including a 99.8\% affirmative on the security investment item. These ceiling responses are best interpreted as conviction rather than as discriminating measurement. Their near-zero variance is also what acquiescence to an affirmatively framed item would produce, so these two items serve as corroboration rather than as load-bearing evidence. That perceived value is strong despite low familiarity is consequential for the adoption argument advanced in Chapter~2. Perception, not demonstrated technical merit, drives adoption decisions~\parencite{Call_Herber_2022, Call_2024}, and in this population high perceived value co-occurs with low awareness rather than following from it.

\subsection{Research Question 3: Anticipated Benefit for Job Performance, Compliance, and Capability}

Research Question~3 concerns belief that Digital Engineering could help in practice, by improving job performance, easing compliance, or strengthening organizational capability. The extent of that belief proves considerable, subject to one qualification. Agreement was substantial on the IT operational items (87.7\%) and the security operational items (82.0\%), and it remained a clear majority on the explicitly compliance-framed items (76.5\%). The compliance grouping was the most reserved, and the single most hedged item in the instrument asked respondents to characterize their own organization's regulatory exposure rather than to endorse a capability in the abstract. That qualification is itself informative, because respondents endorse Digital Engineering's capability to support compliance more readily than they assert that their own organization faces compliance pressures it could address. The result aligns with the compliance imperative documented in Chapter~2. Under that imperative, frameworks such as the NIST Risk Management Framework and the Cybersecurity Maturity Model Certification impose traceability and documentation demands that current practices struggle to sustain~\parencite{Force_2020, DoD_CMMC_2024}. One boundary attaches to this answer and to the previous one together. Research Questions~2 and~3 were answered from a shared pool of items, and the construct-validity analysis found the value dimensions to be a single appraisal rather than separable ones. Both conclusions are therefore better read as one finding at two levels of specificity, and the conceptual distinction between general value and personal applicability stands as untested rather than as confirmed.

These answers also address the four propositions advanced in Chapter~1. Proposition~P1 (awareness of Digital Engineering as a discipline is low) is supported by the familiarity and exposure data. Proposition~P2 (perceived value is nonetheless positive once capabilities are described in functional terms) is supported by the near-universal capability endorsement and the high value ratings. Chapter~4 nonetheless marks it as the most weakly tested of the four, because the instrument's positively framed capability descriptions made disconfirmation unlikely by design.

Proposition~P3 (awareness and value vary by professional field and experience) is supported on both components. The variation by professional field is large and stable across the statistical controls the design permits. Chapter~4 marks its magnitude as provisional, however, given the recruitment frame that cannot be tested in these data. The variation by experience is a genuine experience-group difference rather than a byproduct of cohort composition. The mid-career cohort at 11 to 15 years of experience scores significantly lower on both value composites even after professional field and familiarity are held constant (regression coefficients of $-0.40$ and $-0.53$, both $p<.001$). Because professional field and experience are only weakly associated in these data, the mid-career reservation cannot be explained away by which professionals populate the cohort.

Proposition~P4 (greater familiarity is associated with higher perceived value) is supported by the regression and rank-order analyses (standardized $\beta \approx 0.33$--$0.34$ across the two value composites). These are correlational endorsements from a cross-sectional design, so ``supported'' here denotes consistency with the evidence rather than confirmation.

\section{Interpretation and Implications}

At its center, this dissertation contributes a first empirical baseline of awareness and perceived value of Digital Engineering among Information Technology and Information Assurance professionals. The baseline draws on 1{,}411 respondents in a domain that previously held only theoretical proposition. That contribution is deliberately descriptive. It establishes professional receptiveness; it does not establish that Digital Engineering transfers to these domains or that it delivers value once implemented, a boundary the limitations preserve. Three findings carry implications beyond the immediate research questions.

The largest effect in the data separates the two focal fields. Information Technology respondents outscore Information Assurance respondents by a wide margin, with effect sizes between $d=1.64$ and $d=1.94$ across all three composites. One of those three figures carries a qualification the others do not. Because the Awareness composite blends the low-scoring familiarity item with the six capability-recognition items, its $d=1.64$ is better read as a gap in capability endorsement than as a gap in awareness, and Chapter~4 reports it on that basis. This gap runs counter to the dissertation's expectation that the Information Assurance community would be the most receptive audience. Chapter~2's wider expectation fared better once separated into its parts. The attributes that depend on contact with the methods behaved as predicted, since exposure and familiarity are both low, while relative advantage did not, since perceived value is high across every block. Adoption theory anticipates that pairing, and the present data locate this population precisely there: favorably disposed toward an advantage they have had little opportunity to observe.

Information Assurance respondents reported markedly lower familiarity, with a Q1.1 mean of 1.20 against 2.88 for IT, and lower perceived value. Three readings of the gap remain open. The first is a genuine awareness deficit specific to the Information Assurance community. The second is a valid relevance judgment: the group is heterogeneous, and its operationally oriented members may reasonably conclude that a model-centric and document-centric paradigm fits hands-on defensive work less directly. The third is an artifact of sample composition. Cross-sectional perception data cannot resolve which of these holds.

The focused breakout analysis in Chapter~4 nonetheless narrows the question by ruling out one alternative explanation. The within-domain models confirm directly that the field gap is not an experience effect in disguise: Information Technology exceeds Information Assurance within each of the three interpretable experience cohorts. The candidate explanations that remain carry different imperatives. If the deficit is genuine, the community that the compliance imperative most directly affects is also the least aware of the methods that could address it. That is precisely the condition under which professional-development and curricular initiatives are warranted. If the gap instead arises partly from the recruitment frame, from differential item functioning, or from self-categorization, then replication with representative sampling of the Information Assurance community becomes the priority. Replication would ensure that a measured rather than an apparent deficit directs intervention resources.

A tension left open in Chapter~2 also finds an empirical answer here. The adoption literature holds that perception governs adoption and that preventative innovations such as Model-Based Systems Engineering suffer from low observability that depresses perceived advantage~\parencite{Call_2024}. Chapter~4's results show that when capabilities are described in functional terms tied to recognized professional problems, perceived value is high even absent familiarity. The evidence is therefore consistent with awareness and exposure, rather than skepticism about value, being the more likely constraint on adoption in these domains, among the constraints the instrument measured. Because the design is cross-sectional and correlational, this remains a diagnosis to be tested rather than a causal conclusion the present data can establish; identifying the binding constraint would require the intervention and longitudinal studies proposed below. One step in that inference is assumed rather than tested. The adoption literature's construct is the valuation a professional brings to a real adoption decision, while this study measured stated agreement with described capabilities, most of it from respondents who had never encountered the methods in practice. Whether the second predicts the first is an empirical question these data do not settle, and the intervention studies proposed below would test it.

Two cautions attach to reading the familiarity--value gradient as evidence that exposure reveals value. The gradient is equally consistent with exposure lowering appraisal, since implementation cost and complexity become salient only once professionals engage the methods. The high-familiarity stratum is also self-selected. It contains those who encountered Digital Engineering and stayed; those who tried it and judged it not worth the cost are absent from the sample. Only the proposed intervention and longitudinal studies can therefore establish what direction genuine exposure would produce in the naive majority. If that reading is borne out, it identifies the form effective intervention should take: education and trial opportunity rather than persuasion about benefit.

Read through the boundary-object lens of Chapter~2, the IT-versus-Information-Assurance divide is the predictable signature of a pragmatic knowledge boundary~\parencite{Star_Griesemer_1989, Carlile_2002, Carlile_2004}. The shared, authoritative models of Digital Engineering are positioned to cross that boundary only once awareness brings them within these communities' reach. The high perceived value the study records speaks directly to the crossing condition that \citeauthor{Carlile_2004} identifies. Crossing a pragmatic boundary requires parties to accept the cost of transforming their working knowledge, and that acceptance presupposes that they judge the artifact carrying the transformation to be worth the cost. These professionals already rate the described capabilities as highly valuable. That rating is consistent with, though not equivalent to, a willingness to bear the transformation cost, which only implementation research could confirm. On this reading the transformation would await exposure to the models rather than assent to their worth. Among the constraints this study measured, awareness and exposure are therefore where the binding one sits, as the theory predicts once value is granted. The other costs of crossing a pragmatic boundary, in tooling, in skills, and in organizational commitment, lie outside the instrument, and only the proposed intervention studies can weigh them against awareness.

The bias diagnostics, finally, qualify but do not overturn the value findings. The priming diagnostic argues against a question-framing explanation that would operate uniformly across familiarity, because high-familiarity and low-familiarity respondents diverged sharply rather than converging. It cannot, however, exclude a framing or self-selection effect that itself co-varies with familiarity, since the high-familiarity stratum is self-selected. The wave analysis, by contrast, found no material difference between early and late respondents on any composite ($|d| \leq 0.04$, all Mann--Whitney $p > .5$), so response timing yields no evidence of self-selection along the early-to-late dimension. Reported together, these diagnostics support a measured conclusion. Perceived value is genuine and substantial. The defense-and-aerospace recruitment frame examined in the limitations still counsels reading the favorable estimates as upper-leaning for the full target population.

These results contribute the empirical foundation that Chapter~1 identified as missing. They establish the baseline that any subsequent implementation, case-study, or curricular research would draw on. The findings are consistent with the theoretical premise that Digital Engineering's capabilities resonate with the documented problems of IT and Information Assurance practice, and they show that low awareness coexists with high perceived value in this population. Between those two candidate constraints, awareness and exposure are the more likely brake on adoption. No other barrier entered the instrument, so this ranking holds within the pair the study assessed and says nothing about implementation cost, tooling maturity, skills availability, or organizational readiness. The work thereby completes the second step in the logical pathway articulated in Chapter~1, which runs from demonstrated value in defense and aerospace, through established professional receptiveness, toward future investigation of practical implementation.

\section{Research Significance Revisited}

Chapter~1 framed the significance of this research along four dimensions. The results now permit reconsidering each against evidence rather than expectation.

\subsection{Academic Significance}

This work supplies empirical evidence to a domain previously characterized by theoretical proposition rather than data, realizing the academic significance Chapter~1 anticipated. The literature reviewed in Chapter~2 contained no investigation of whether IT and Information Assurance professionals perceive value in Digital Engineering capabilities. The systematic mapping of the field by \citeauthor{Wooley_Womack_2025} likewise found no research addressing enterprise IT or Information Assurance applications~\parencite{Wooley_Womack_2025}. In method, the study extends the diffusion-of-innovations perception surveys of \citeauthor{Call_2024} and \citeauthor{Henderson_McDermott_Salado_2024} from the systems engineering community to the Information Technology and Information Assurance population~\parencite{Call_2024, Henderson_McDermott_Salado_2024}. That extension is the point rather than an incidental feature, because this population is the precise locus of the transferability question the dissertation raises, and no prior survey had measured that population. Present findings supply the missing perception data and, in doing so, test the theoretical framework that connects Digital Engineering capabilities to documented IT and Information Assurance challenges. On the value dimension that framework survives, since practitioners recognize substantial value, while the awareness dimension reframes the research agenda toward education and exposure.

Two methodological results of independent interest follow as well. One is measurement-theoretic. That low familiarity coexists with high capability-endorsement shows empirically why familiarity with a discipline and recognition of its capabilities require separate measurement, a distinction that the layered Section~1 and Section~2 design was built to capture. A third construct, verified knowledge, requires a third measure that this instrument did not carry. Another arose as an unplanned by-product: the anonymous open-link mode of data collection required defending the data against the automated, coordinated, and synthetic responding that increasingly threatens online surveys~\parencite{Pozzar_Bots_2020, Westwood_LLM_2025, Jaffe_NotBots_2026}. Mounting that defense required developing, validating, and openly releasing a duplicate-vector response-authenticity screen paired with human adjudication of borderline cases (Appendix~B). Of the 71 identical vectors submitted in the one coordinated burst that reached the instrument, the pre-registered timing screen removed the 70 that arrived too quickly and the response-vector rule removed the single submission whose slower timing would otherwise have survived. The study offers the artifact as a reusable, auditable contribution and develops it more fully in a companion methods manuscript.

\subsection{Industry Significance}

Findings deliver an actionable diagnosis of the kind Chapter~1 envisioned for industry: a characterization of current awareness and perceived value that can inform adoption strategy across the professional community. Organizations contemplating Digital Engineering adoption face a receptive rather than a skeptical audience, since perceived value and recommendation intent are high. Receptiveness is a fact about the audience rather than about the technology. These results speak to what professionals expect Digital Engineering to do for them and leave what it does in deployment to the implementation research the future-work agenda specifies. They must, however, contend with a pervasive awareness deficit that places these capabilities outside the consideration set during tool selection and strategic planning. The breakout analysis sharpens the guidance further by identifying where receptiveness is weakest. The Information Assurance community, which the compliance landscape most directly affects, reports the lowest awareness and value. If that deficit proves genuine rather than an artifact of sampling, it marks this community as the segment requiring the most deliberate engagement. Vendors, service providers, and standards bodies can therefore provisionally concentrate communication and enablement on the audiences that the data identify as least served, principally the Information Assurance community, rather than allocating effort uniformly.

What form that effort should take depends on how the deficit is read. If Information Assurance respondents report lower value because Digital Engineering has not reached them, then awareness-building and trial opportunity are the appropriate instruments. If their lower ratings instead reflect calibrated professional skepticism toward affirmatively framed claims about security capability, as Chapter~4 raises among the candidate explanations, then the same audience calls for demonstrated evidence of benefit in security settings rather than broader exposure. The present data cannot separate the two, and the distinction governs the design of an intervention rather than the standing of the finding.

The instrument does not, however, license capability-by-capability prioritization. Because many items pair a capability with an asserted benefit, and some bundle more than one function, the data speak more securely to which audiences are least reached than to which individual pillar is most or least valued. The instrument refinements below therefore include rewording each item to a single clause so that per-pillar demand reads cleanly.

\subsection{National Significance}

Respondents endorsed, at high rates, the propositions that Digital Engineering could improve compliance with security frameworks and could shorten vulnerability identification and remediation. These are precisely the capabilities that the expanding requirements of SP~800-171 and the Cybersecurity Maturity Model Certification demand~\parencite{NIST_SP_800_171_R3_2024, DoD_CMMC_2024}. These results bear directly on the national dimension of significance, which concerns national security and the protection of societal infrastructure. Chapter~1 argued that authoritative traceability of the kind Digital Engineering provides could reduce the burden of compliance verification for operators of federal and national security systems while improving the currency of compliance evidence. Reservations voiced on organization-specific compliance items indicate that demonstrating concrete compliance outcomes will advance the case for Digital Engineering in this setting more effectively than asserting general capability. The future-research agenda below takes up this point directly. These findings address only indirectly the dependency-mapping problem that opened this dissertation: the difficulty of tracing the cascading first-, second-, and third-order effects of a vulnerability across interconnected systems. Survey items measured perceived value for vulnerability identification and remediation and for authoritative traceability in general, not for cascading-impact assessment as a distinct use case. The study therefore establishes professional receptiveness to the class of capability to which dependency mapping belongs without measuring that specific application on its own.

\subsection{Societal Significance}

Compliance documentation consumes scarce capacity in organizations that operate under tight resource constraints: community health centers, tribal organizations, under-resourced school districts, small municipalities, and social-service non-profits. Chapter~1 located the societal significance of this work in exactly these settings. The argument advanced there was distributional. Because the manual labor and specialized expertise that compliance demands fall hardest where capacity is smallest, a capability that lowers that threshold would benefit these organizations most. Such a capability could narrow rather than widen the gap between the security and compliance posture available to well-funded institutions and that available to the smallest and least-resourced organizations.

Sampling did not specifically reach those organizations, so the study cannot speak to adoption within them. The study measured receptiveness in a comparatively well-resourced, defense-and-aerospace-weighted population, as the limitations detail, so it establishes the precondition for the accessible population rather than demonstrating it for the under-resourced organizations the argument concerns.

What it does establish is the precondition on which that broader benefit depends. Across a diverse professional population, respondents perceived Digital Engineering's capabilities as valuable, at agreement rates above 80\%, for precisely the documentation, traceability, and verification work that burdens resource-constrained organizations most heavily. That perceived value held even among the large majority of respondents unfamiliar with the umbrella discipline. By confirming that broad professional receptiveness, the study supplies the empirical warrant for the next step: designing and evaluating accessible, lower-cost Digital Engineering implementations suited to organizations with limited specialized expertise. That work would pursue the equity benefit anticipated in Chapter~1 deliberately rather than leave it to chance.

Nothing about the next step is automatic. Digital Engineering as currently practiced is itself expertise- and tooling-intensive, so extending it to the least-resourced organizations depends on accessible, lower-cost implementations that do not yet exist. This chapter therefore frames the distributional benefit as a precondition and a design-science agenda rather than as a demonstrated outcome.

\section{Limitations}

Several limitations bound the conclusions above, some anticipated in Chapter~3 and some surfaced by the results.

Non-probability convenience sampling limits generalizability to the broader population of IT and Information Assurance professionals, so the results should be read as indicative of the accessible population rather than as definitive measures of the entire professional community. Self-selection bias may have over-represented professionals already favorably disposed toward Digital Engineering. The wave analysis did not corroborate this specific concern: early and late respondents did not differ materially on any composite ($|d| \leq 0.04$), which weighs against response-timing self-selection. That analysis cannot, however, exclude a self-selection process operating uniformly across the collection window. Social desirability bias may also have inflated value responses, particularly given the affirmatively framed capability statements, although the anonymous design was intended to mitigate this.

A specific composition concern attends the recruitment strategy. Recruitment relied principally on outreach through the candidate's LinkedIn network, which is itself concentrated in the defense, aerospace, and federal sectors, the very settings in which Digital Engineering originated and is most actively practiced. No industry filter was applied to the invitation, so the frame's composition rather than a selection rule produced the sector weighting. That weighting still tilts the realized sample toward professionals more likely than the broader IT and Information Assurance population to have encountered Digital Engineering. Such weighting plausibly elevates both the awareness and the perceived-value estimates above what a sector-balanced sample would yield. Paid promotion widened circulation beyond that network during the later weeks of collection, and its effect on composition runs in both directions. Reaching professionals outside the candidate's contacts works against the network skew, while delivery through an advertising algorithm the study neither controlled nor recorded introduces a selection mechanism of its own. Wave analysis, however, found that the later respondents it attracted did not differ materially from earlier ones on the value composites. Read conservatively, then, the value findings are upper-leaning estimates for a defense-and-aerospace-weighted accessible population rather than point estimates for the entire professional community.

This limitation is the most consequential the study carries, because the instrument recorded professional field and experience but no measure of employment sector or organization type. No direct adjustment for the frame effect is therefore possible from these data. The proxy adjustments reported below, which hold recent exposure and familiarity-stratum composition constant, narrow the concern without resolving it. Future replication should carry a sector item so that this confound can be modeled rather than reasoned about. The comparative findings, which contrast groups drawn from the same sample, are more robust to this composition than the absolute levels, because a shared recruitment bias affects the groups being compared in common.

Several limitations follow from the realized data. The instrument produced four professional-field categories rather than the planned dichotomy, and the primary comparison contrasted Information Technology with Information Assurance while retaining Engineering and Other for descriptive context only. A different operationalization of field would therefore change the comparative section. The two investment-willingness items exhibited ceiling effects that rendered them near-invariant and therefore uninformative for subgroup contrasts. Several stems also bundle more than one capability or outcome within a single statement, among them the vulnerability item that pairs identification with remediation, so agreement on those items cannot be attributed to one clause of the stem. Chapter~4 marks the affected items where it reads them individually.

The 16+ year experience cohort, at only eighteen respondents, was too small to characterize, so the experience analysis rests on the three larger cohorts. Composition of that variable also bounds what population the study reached, and it does so on evidence the instrument measured directly. Respondents with more than fifteen years in the field number eighteen, or 1.3\% of the sample, while 40.1\% report five years or fewer. A cross-section of the IT and Information Assurance professions would carry a substantially larger senior segment than that. Realized this way, the distribution indicates that the recruitment channels reached a predominantly early-career population rather than the professions at large, which is the clearest evidence available here about the accessible frame the study actually obtained. It also reinforces the recruitment-frame caution attached to the field comparison below, since channels that skew this sharply on tenure may skew on other unmeasured attributes. The cross-sectional design captures perceptions at a single point in time and cannot speak to how awareness or perceived value would change following exposure or education.

A further caution attends the headline field comparison. The Information Technology versus Information Assurance effect sizes ($d=1.64$ to $1.94$, contrasting 501 IT and 501 Information Assurance respondents) are exceptionally large by the standards of survey research, and an effect of that magnitude between two adjacent professional communities itself invites scrutiny.

Beyond the experience confound that the breakout analysis excludes, four alternative explanations remain open. One account is a recruitment-frame difference in how the two groups were reached; another is differential item functioning in how security professionals calibrate affirmatively framed statements. The gap could also reflect a self-categorization effect in the choice of professional label: the combined Information Assurance/Cyber Security response option spanned both governance-oriented and operationally oriented professionals and so mixed two roles the instrument cannot separate. A fourth possibility is that the lower Information Assurance ratings represent calibrated professional judgment, with respondents discounting affirmatively framed claims about security capability because their expertise equips them to appraise those claims critically. That reading gains some support from the observation that Information Assurance respondents rate the security value of Digital Engineering below the rating Information Technology respondents assign it, and it would leave the measured difference intact while changing what the difference means. Present data cannot adjudicate among these, and the magnitude should not be taken at face value as a pure awareness difference.

Three analyses nonetheless bear on the question. Chapter~4's differential-item-functioning screen weakens the measurement-artifact account, and adding a recent-exposure indicator to the regression leaves the field coefficient essentially unchanged. Holding constant the composition of the bimodal familiarity distribution, which divides the sample into a large naive stratum and a small well-acquainted one, likewise retains 87\% of the field coefficient.

The anonymous instrument recorded no employment sector and no recruitment channel, so which channel reached which community cannot be reconstructed. An effect this large is therefore as consistent with an unmeasured frame difference as with a true awareness gap, and confirming it is a task for replication rather than for reinterpretation of these data. The appropriate reading is that the field difference is a strong but provisional result. Establishing its size in the population requires a probability or quota sample that records recruitment channel and employment sector, a design the future-research agenda specifies.

A related boundary bears on the adoption argument. The instrument measured awareness, exposure, and perceived value, and it measured no other candidate barrier to adoption. Several plausible barriers went unmeasured. The instrument asked nothing about implementation and licensing cost, about tooling maturity, about the availability of modeling skills, about competing organizational priorities, or about an organization's readiness to absorb a new method. It also left aside the effort expectancy, facilitating conditions, and social influence that the Unified Theory treats as determinants alongside performance expectancy~\parencite{Venkatesh_2003}. Comparing awareness against value perception is therefore the most this design supports. Where the study finds awareness low and value high, the reasonable inference is that awareness and exposure constrain adoption more than doubt about benefit does. Which barrier binds hardest overall remains open, and a barrier inventory administered alongside the awareness and value measures would settle it.

Most fundamentally, the study measured perceived value and anticipated benefit rather than actual experienced benefit. Positive perceptions establish receptiveness, yet they do not demonstrate that Digital Engineering would deliver value if implemented, nor do they establish organizational readiness for adoption. The value finding is further bounded by the instrument's construction. All thirty items are affirmatively keyed, and the instrument carried no reverse-coded or knowledge-verification items, so a respondent disposed to agree encountered no item on which agreement would register as disconfirmation. Near-universal endorsement is therefore close to the strongest result the instrument could produce, and \emph{receptiveness}, rather than demonstrated conviction, is the ceiling of what the value data support. The negligible straight-lining rate and the priming diagnostic bound crude and uniform acquiescence, but neither can exclude a mild yea-saying tilt that co-varies with familiarity.

Because the design is anonymous, professional field, the basis of the central comparison, is self-reported and cannot be verified; anonymity trades individual verifiability for candor. The response-authenticity screen detects coordinated and duplicated submissions, not individual misrepresentation of role, so the guard against misclassification is indirect. The field groups exhibit coherent, theory-consistent internal structure rather than noise, straight-lining is negligible, and an unpaid anonymous survey offers little evident incentive to misrepresent one's field. The single coordinated attack that reached the instrument was a duplication burst that the screen removed. Individual misrepresentation of role is, by design, undetectable by any content-based screen, so this residual is acknowledged rather than claimed away. A design that records verifiable role or employment, or a screened professional panel, would close this gap and is specified in the future-research agenda.

The supplementary post-hoc analyses reported in Chapter~4 refine two of these cautions and introduce a third. The differential-item-functioning screen found only two of the twenty-three composite items functioning differently across the two professional fields, both by small margins, which weakens the measurement-artifact account of the field gap without excluding it entirely. The construct-validity assessment found the three composites to be reliable but strongly intercorrelated rather than discriminant, so this dissertation treats them as summative indices of related perceptions rather than as independent constructs. Any claim that awareness and perceived value are separable dimensions is correspondingly bounded. Because Research Questions~2 and~3 draw on the same item pool, their separate answers do not constitute independent findings. Because the study added these analyses after data collection and did not specify them in the approved proposal, they are exploratory and should be read as robustness corroboration of the reported findings rather than as pre-registered confirmatory tests.

\section{Future Research}

Five complementary lines of future research follow from these findings.

\subsection{From Perceived to Measured Value}

Case-study and pilot-implementation research in enterprise IT and Information Assurance settings offers the most direct extension, converting perceived value into measured value. Such work would address the measurement deficit that Chapter~2 documented across the broader Digital Engineering literature, where most claimed benefits rest on perceived rather than measured evidence~\parencite{Henderson_Salado_2021, Rogers_Mitchell_2021}. A pilot would test whether the value this population perceives materializes in practice, under what organizational conditions, and at what cost. Coupling such a pilot to an explicit benefits-measurement framework of the kind \citeauthor{Henderson_2023_Metrics} developed would convert anecdotal outcomes into quantified evidence suitable for enterprise business cases~\parencite{Henderson_2023_Metrics}.

\subsection{Explaining the Field and Experience Patterns}

Why Information Technology so consistently outpaces Information Assurance calls for investigation in its own right. Because cross-sectional perception data cannot distinguish a genuine awareness deficit from a sampling artifact, future work should use stratified or probability sampling within the Information Assurance community specifically. Qualitative inquiry should accompany that sampling, probing why these professionals report lower familiarity and value despite operating under the compliance pressures that Digital Engineering is positioned to address.

The breakout analysis also surfaced a pattern that merits dedicated study: the genuine mid-career reservation at 11 to 15 years of experience. The reservation persists across both focal fields and net of statistical control for field and familiarity, yet it lies outside what the main-effects model explains and remains substantively unaccounted for. Three hypotheses, each grounded in existing theory, are worth testing. The first is a sunk-cost or switching-cost account: mid-career professionals have the deepest personal investment in the incumbent document-centric methods and tooling. The perceived cost of a paradigm shift is therefore highest for them and depresses their appraisal of its value. Earlier-career professionals carry less such investment, and senior professionals evaluate the shift as a portfolio choice rather than a personal retooling. The second is a diffusion-of-innovations account~\parencite{Call_2024}: the 11-to-15-year cohort may occupy a pragmatist position on the adoption curve, past the exploratory openness of newer entrants but short of the strategic vantage of senior staff. Pragmatists are precisely the segment that a preventative innovation with low observability persuades least. The third is a role-structure account: at this tenure professionals often hold hands-on senior-technical or first-line-management roles in which they would absorb the operational burden of adoption without yet holding the authority to direct it or the budget to resource it. That position rationally tempers enthusiasm.

These accounts are not mutually exclusive, and they make different predictions about which interventions would move the cohort, so distinguishing them matters practically. Interviews and focus groups would be well suited to the task, since professional narrative captures a mid-career shift in how professionals appraise emerging methods more naturally than additional survey items would.

\subsection{Intervention, Longitudinal, and Replication Studies}

Because the evidence is consistent with awareness, rather than perceived value, being the limiting factor for adoption, intervention research becomes a natural next step and the direct test of that reading. Studies that measure awareness, perceived value, and adoption intention before and after structured exposure or training would test whether closing the awareness gap shifts adoption intent. They would also connect this work to the enterprise adoption frameworks synthesized in Chapter~2~\parencite{McDermott_2024}. Longitudinal designs would address the temporal limitation of the present cross-sectional snapshot by tracking how perceptions evolve as Digital Engineering diffuses into enterprise contexts. Replication against probability frames, across additional sectors, and in other national contexts would establish the external validity that convenience sampling cannot, and it would test whether the field and experience patterns observed here generalize. Any of these designs should carry a barrier inventory, measuring cost, tooling maturity, skills availability, and organizational readiness alongside awareness and perceived value, so that the obstacles to adoption can be ranked against one another rather than assessed one at a time.

\subsection{Instrument and Methodological Refinement}

The instrument itself warrants refinement. Its ceiling effects on the investment-willingness items indicate that future versions should replace binary or ternary willingness questions with graduated or scenario-anchored measures capable of discriminating among the favorably disposed. That low familiarity coexists with near-universal agreement also argues for incorporating reverse-coded items and knowledge-verification items carrying scoreable correct answers. Such items would measure comprehension directly, which the present design could only approach through recognition, and they would separate substantive knowledge from acquiescence. Refining the field item to elicit the four categories deliberately, rather than recovering them after the fact, would support a planned four-group comparison in place of the proposal's dichotomy.

The measurement analyses in Chapter~4 point to two further, more specific refinements. First, the average variance extracted for each composite fell just below the conventional 0.50 benchmark, which reflects item sets that each span a heterogeneous range of capabilities rather than a single tightly defined facet. Tightening each composite toward a more homogeneous set of indicators, or lengthening it with items that load more cleanly on one construct, would raise convergent validity in a future administration. Second, the heterotrait--monotrait analysis showed the IT-value and security-value dimensions to be empirically inseparable. A revised instrument should therefore measure perceived value as a single scale, which the present data support. If theory instead predicts a substantive IT-versus-security distinction, the revision should redesign the items so that the two facets are worded to diverge rather than to restate a common appraisal in two settings. Several of these refinements---including reverse-coded, knowledge-verification, and per-capability items---are carried forward into the reproducibility artifact described under Dissemination and Reproducibility below. A subsequent study can therefore adopt them without reconstructing the instrument from the proposal.

\subsection{Technical and Design-Science Directions}

Technical barriers identified in Chapter~2 define a final, design-oriented agenda. Developing and validating modeling-language extensions and domain-specific profiles for representing IT services, network configurations, and security controls would close the gap between the general-purpose modeling languages of systems engineering and the specific representational needs of enterprise IT and Information Assurance. Design-science research could build and evaluate a Digital Engineering reference implementation that maintains an authoritative compliance traceability thread across the Risk Management Framework and the Cybersecurity Maturity Model Certification, then measure the reduction in audit-preparation effort against current document-centric practice. Such research would extend the preliminary reference-model work that motivated this study~\parencite{Bonar_Hastings_2024, CMU_SEI_DevSecOps_2023}.

A distinct and equity-oriented strand of this agenda would target the small and under-resourced organizations identified earlier. Design-science work could develop and evaluate lightweight, low-cost, and template-driven Digital Engineering implementations that minimize the specialized expertise required for adoption. These implementations would be packaged with pre-built compliance profiles for the regulations these organizations most commonly face, such as HIPAA for community health providers and FERPA for under-resourced schools. Field studies in such settings, building on the precedent of Model-Based Systems Engineering in rural healthcare~\parencite{Berg_2023}, would then measure whether the documentation and compliance burden actually falls and whether organizations genuinely redirect scarce staff capacity toward direct service provision. This line would convert the distributional promise identified here into measured outcomes for the organizations that stand to benefit most.

Related work on tool interoperability would address the integration challenge of connecting systems-engineering modeling environments to the IT service management, security information, and configuration platforms that enterprises already operate. Finally, Chapter~2 documented that inadequately secured digital twins expand the attack surface they are meant to reduce. Research into securing enterprise digital twins within an integrated Digital Engineering framework would therefore ensure that adoption advances security posture rather than undermining it~\parencite{Suhail_Iqbal_Jurdak_2024}.

\subsection{Dissemination and Reproducibility}

Two components of this work are positioned for peer-reviewed dissemination in their own right. The first is the structured cross-disciplinary review of Chapter~2, which delineates a research gap that no prior synthesis has mapped. The second is the perception survey of Chapters~3 and~4, whose sample of 1{,}411 professionals is large by the standards of this research area. To support replication and reuse, the survey instrument, the scripted analysis pipeline, and a synthetic dataset that reproduces the reported results are openly available in a public research repository at \url{https://doi.org/10.17605/OSF.IO/VKDTE}. This privacy-preserving reproduction dataset is distinct from the separate, purpose-built synthetic ground-truth set used to establish the operating characteristics of the response-authenticity screen; the two serve different roles and are not conflated. Because the instrument was fully anonymous and collected no identifiers, the de-identified response data underlying the study are available from the author on request. Their public deposit is under consideration, subject to the anonymity and exempt-review terms of the study.

\section{Concluding Remarks}

This research set out to determine whether IT and Information Assurance professionals recognize the potential of Digital Engineering for their work, and it answers that question with empirical data where previously only proposition existed. Respondents are largely unfamiliar with Digital Engineering as a discipline, yet they endorse its constituent capabilities and perceive substantial value in them for IT operations, security operations, and compliance. That perceived value is strong enough that nearly all respondents would recommend their organizations explore adoption. Of the two barriers this work measured, awareness and exposure are the binding one, and skepticism about value is not. Adoption faces other obstacles the study did not assess, among them implementation cost, tooling maturity, and organizational readiness, so this finding locates one constraint rather than ranking the full set. That deficit appears most pronounced within the Information Assurance community that the expanding compliance landscape most directly affects. This reading remains associational, and the magnitude of the gap warrants the cautions raised in the limitations.

Chapter~2's theoretical lens frames what such receptiveness makes possible. Should the models of Digital Engineering prove able to act as boundary objects, carrying transformed knowledge across the pragmatic divide that separates systems engineering from Information Assurance and IT service management, then the willingness measured here becomes the precondition for that transfer. By establishing that perceived value is high while awareness is low, this dissertation provides the evidentiary foundation for the next stage of inquiry: demonstrating, measuring, and enabling the practical application of Digital Engineering capabilities in enterprise IT and Information Assurance. The need in those domains is documented, and the receptiveness, this study shows, is already present.
